% Vietnamese translation for OI Public Library
% Source-File: IOI中国国家候选队论文/2008/Day2/1.程芃祺《计算几何中的二分思想》/计算几何中的二分思想.doc
\documentclass[11pt]{article}
\usepackage{oipl-vn}

\newcommand{\dist}{\operatorname{dist}}

\title{Tư tưởng chia đôi trong hình học tính toán}
\author{Cheng Pengqi\\
Trường Trung học số 1 Quý Dương}
\date{2008}

\begin{document}
\maketitle

\origtitle{Binary-search thinking in computational geometry}
\sourcepath{IOI China candidate team papers / 2008 / Day 2 / Cheng Pengqi}

\begin{abstract}
Bài viết trình bày ngắn gọn tư tưởng chia đôi trong hình học tính toán, rồi
thông qua các bài toán ví dụ để minh họa cách áp dụng tư tưởng ấy, qua đó thể
hiện ưu thế gọn gàng và hiệu quả của chia đôi trong hình học tính toán.
\end{abstract}

\noindent\term{Từ khóa:} hình học tính toán; chia đôi.

\section{Dẫn Nhập}

Tư tưởng chia đôi đã có từ xưa. Thiệu Tử từng nói: ``Một chia làm hai, hai
chia làm bốn, bốn chia làm tám.'' Chính dựa trên cách nghĩ ấy, tổ tiên chúng
ta đã sáng tạo ra Thái Cực và Bát Quái. Trong thời đại thông tin ngày nay,
trí tuệ cổ xưa này vẫn tỏa sáng; nó thấm vào nhiều ngành học mới và được phát
huy mạnh mẽ, trong đó có hình học tính toán. Trên nền tảng ấy đã sinh ra vô
số thuật toán kinh điển, chẳng hạn dùng chia để trị để tìm cặp điểm gần nhất,
bao lồi, tam giác hóa, và cây nhị phân phân hoạch không gian.

Trong các kỳ thi tin học những năm gần đây, ngày càng xuất hiện nhiều bài toán
về hình học tính toán. Trong đó, khá nhiều bài phức tạp có thể được giải quyết
đơn giản nhờ tư tưởng chia đôi. Nắm được tư tưởng này rõ ràng là có thêm một
công cụ sắc bén để xử lý các bài toán liên quan.

\section{Phân Tích Ví Dụ}

Dưới đây, ta cùng khảo sát ứng dụng của tư tưởng chia đôi qua một số bài toán.

\subsection{Ví dụ 1: \textit{Simplified GSM Network} (2005 ACM/ICPC World Finals)}

\paragraph{Ý nghĩa bài toán.}
Cho tọa độ của \(B\) trạm tín hiệu \((1\le B\le 50)\) và \(C\) thành phố
\((1\le C\le 50)\), với trị tuyệt đối của mỗi tọa độ không vượt quá \(1000\).
Mỗi thành phố sử dụng trạm tín hiệu gần nhất. Cho \(R\) mô tả tuyến đường nối
giữa các thành phố \((1\le R\le 250)\) và \(Q\) truy vấn \((1\le Q\le 10)\).
Hãy tìm số lần chuyển trạm tín hiệu ít nhất khi liên lạc giữa hai thành phố
tương ứng.

\paragraph{Phân tích.}
Rõ ràng bài toán cần tìm đường đi ngắn nhất; điểm then chốt nằm ở cách tính
trọng số của từng tuyến. Đề bài cho biết mỗi thành phố dùng trạm tín hiệu gần
nhất, tức vị trí của thành phố nằm trong quỹ tích các điểm trên mặt phẳng gần
trạm ấy nhất. Quỹ tích các điểm gần một điểm nhất chính là biểu đồ Voronoi.
Vì vậy, trọng số của mỗi tuyến bằng số cạnh Voronoi mà đoạn thẳng ấy cắt qua.
Từ đây có thể thấy bài toán là một tổ hợp đơn giản của biểu đồ Voronoi và
đường đi ngắn nhất; chỉ cần giải hai phần riêng rẽ.

Phân tích trên dẫn tới cách trực tiếp nhất: trước hết dựng biểu đồ Voronoi,
sau đó chạy thuật toán đường đi ngắn nhất. Độ phức tạp thời gian của cách này
đủ để qua bộ kiểm thử, nhưng việc lập trình dựng biểu đồ Voronoi rất phức tạp
và khó gỡ lỗi; trong môi trường thi thật, tính thực dụng không cao. Có cách
nào tốt hơn không?

Xét một đoạn thẳng \(l\). Nếu hai đầu mút của \(l\) thuộc cùng một trạm tín
hiệu thì hiển nhiên trọng số của \(l\) bằng \(0\). Ngược lại, chia \(l\) tại
trung điểm thành hai phần \(l_1\) và \(l_2\) (điểm chia không nằm trên cạnh
Voronoi). Khi đó trọng số ứng với \(l\) đương nhiên bằng tổng trọng số của
\(l_1\) và \(l_2\).

Tiếp tục thao tác tương tự với \(l_1\) và \(l_2\). Vì không thể chia mãi mãi,
khi khoảng cách giữa hai đầu mút nhỏ hơn một hằng số dương nhỏ \(\varepsilon\)
nào đó (thực nghiệm cho thấy \(\varepsilon\in[10^{-10},10^{-5}]\) đều được
chấp nhận), đồng thời hai đầu mút thuộc các trạm tín hiệu khác nhau, thì giữa
hai điểm ấy nhất định tồn tại một cạnh Voronoi cắt đoạn thẳng; trọng số của
đoạn này là \(1\).

Dùng cách trên, ta có thể tính trọng số của các cạnh một cách đơn giản mà
không cần dựng biểu đồ Voronoi, tránh được mã dài dòng. Sau đó dùng Floyd,
Dijkstra, hoặc một thuật toán đường đi ngắn nhất khác là giải xong toàn bộ
bài toán.

\paragraph{Tiểu kết.}
Hình học tính toán có nhiều mô hình kinh điển như biểu đồ Voronoi. Chúng giúp
ta suy nghĩ thuận tiện, nhưng cũng dễ làm ta mắc vào quán tính tư duy cố hữu.
Khi mô hình ban đầu không còn tiện để giải bài, ta cần mở một lối khác, vượt
ra khỏi cách nghĩ truyền thống.

Trong ví dụ này, tư tưởng chia đôi tách một đoạn thẳng thành hai nửa và dùng
chia để trị để giải bài toán, nhờ đó biến phức tạp thành đơn giản, dễ lập
trình, đồng thời tránh được lối mòn dựng biểu đồ Voronoi.

\subsection{Ví dụ 2: \textit{Collecting Luggage} (2007 ACM/ICPC World Finals)}

\paragraph{Ý nghĩa bài toán.}
Cho tọa độ các đỉnh của một đa giác đơn \(N\) cạnh \((3\le N\le 100)\). Vali
di chuyển từ đỉnh đầu tiên dọc theo biên đa giác với vận tốc \(V_l\). Người
xuất phát từ điểm \((p_x,p_y)\) với vận tốc \(V_p\), trong đó
\(0<V_l<V_p\le 10000\) (đơn vị: mét/phút). Mọi tọa độ đều là số nguyên có trị
tuyệt đối không quá \(10000\) (đơn vị: mét). Hãy tìm thời gian sớm nhất để
người lấy được vali, chính xác tới giây.

\paragraph{Phân tích.}
Thoạt nhìn, bài toán có vẻ liên quan mật thiết tới đường đi ngắn nhất. Tuy
nhiên điểm đích thay đổi theo thời gian, có thể là bất cứ điểm nào trên biên
đa giác; thứ cần tìm là một quá trình động. Tính thời gian cần để tới một đỉnh
là việc đơn giản, nhưng làm sao đến được đích cuối cùng mới là khó khăn của
bài toán.

Trước hết xét trường hợp đơn giản: người đứng ở vị trí ban đầu, vali chuyển
động trên một đường thẳng, ban đầu ở \((x_0,y_0)\), với vận tốc
\(V_l=(a,b)\). Người lấy được vali tại thời điểm \(t\); khi đó tọa độ vali là
\((at+x_0,bt+y_0)\), và ta có quan hệ
\[
  (at+x_0-p_x)^2+(bt+y_0-p_y)^2=(V_p t)^2.
\]
Dùng công thức này, có thể tính khoảng cách ngắn nhất động từ một điểm bất kỳ
tới một đường thẳng; với đoạn thẳng trong đề bài, chỉ cần thêm điều kiện rằng
điểm đích nằm trên đoạn.

Sau khi có công thức khoảng cách, tưởng như phần lớn vấn đề đã được giải
quyết. Nhưng khi đưa nó vào đa giác, ta nhận ra mọi chuyện không đơn giản như
tưởng tượng: một điểm không nhất thiết đi thẳng được tới một cạnh, và cũng
không nhất thiết đi thẳng được tới toàn bộ phần có thể tới của cạnh ấy. Vì
vậy còn phải tìm miền có thể tới, rồi liệt kê các trường hợp người và vali di
chuyển tới từng điểm, từng cạnh; độ phức tạp của vấn đề tăng lên rất nhiều.

Làm thế nào để đơn giản hóa quá trình phức tạp này? Mấu chốt nằm ở việc
``biến động thành tĩnh'', chuyển quá trình động trong đề bài thành bài toán
tĩnh. Giả sử người ``tạm dừng'', còn vali đi trước trong thời gian \(t\) đến
một điểm nào đó. Nếu coi vali đứng yên trên một cạnh, ta chỉ cần thêm một đỉnh
mới và tách cạnh ấy thành hai đoạn, rồi rất dễ tính khoảng cách ngắn nhất.
Sau đó cho người đi trong thời gian \(t\), kiểm tra xem có lấy được vali hay
không, tức so sánh khoảng cách tìm được với quãng đường người đi được.

Cách làm này cho biết người có thể hoàn thành hành động trong thời gian quy
định hay không, nhưng không thể liệt kê từng giây vì hiệu quả quá thấp. Có
thể dùng chia đôi để phán định không? Điều này đòi hỏi tính đơn điệu: nếu
trong thời gian \(t\) vừa đủ đi từ một đỉnh có thể tới đến vị trí vali, thì ít
hơn \(t\) không thể tới, còn nhiều hơn \(t\) chắc chắn có thể tới.

Ta chứng minh trước rằng ít hơn \(t\) thì không tới được. Giả sử
\(t'<t\). Gọi \(d_l\) là khoảng cách đường thẳng giữa hai vị trí của vali ở
hai thời điểm \(t'\) và \(t\), \(L\) là đường đi ngắn nhất tới vị trí vali ở
thời điểm \(t'\), và \(d\) là khoảng cách đường thẳng tương ứng. Từ bất đẳng
thức tam giác, hiệu hai cạnh của tam giác nhỏ hơn cạnh thứ ba, nên
\[
  d>V_p t-d_l.
\]
Lại do \(V_l<V_p\) và \(d_l<V_l(t-t')\), ta có
\[
  L\ge d>V_p t-d_l>V_p t-V_l(t-t')>V_p t-V_p(t-t')=V_p t'.
\]
Vì vậy trong thời gian \(t'\) không thể tới. Chứng minh rằng nhiều hơn \(t\)
thì chắc chắn tới được cũng tương tự, vẫn chỉ cần dùng quan hệ ba cạnh của tam
giác, nên không nhắc lại ở đây.

Quay lại bài gốc, trước hết cần xây dựng đồ thị nhìn thấy được và tính đường
đi ngắn nhất từ điểm xuất phát tới các đỉnh. Vì tính đơn điệu đã được đảm bảo,
ta chỉ cần chia đôi thời gian \(t\), tìm vị trí của vali sau thời gian \(t\),
thêm đỉnh mới rồi tính đường đi ngắn nhất tới các đỉnh, sau đó so sánh để tìm
điểm tới hạn. Một cận rất tốt là từ thời gian nhỏ nhất để tới biên đa giác
đến thời gian lớn nhất, vì cận này bảo đảm phủ toàn bộ phạm vi đa giác; nhưng
điểm đạt cực trị không nhất thiết là đỉnh, còn tính khoảng cách điểm--cạnh khá
phiền. Một cận khác tiện tính hơn là từ \(0\) đến thời gian đi hết tổng của
khoảng cách tới đỉnh gần nhất và nửa chu vi; thực nghiệm cho thấy hiệu quả
khá tốt.

Sau khi hoàn thành phần chính của thuật toán, còn một chi tiết rất quan trọng:
nếu đúng ra cần \(N+0.5\) giây (với \(N\) là số tự nhiên), thì quy tắc làm
tròn xuống ở phần ``bốn'' và làm tròn lên ở phần ``năm'' sẽ khiến số giây ứng
với hai cận luôn khác nhau, làm chương trình rơi vào vòng lặp vô hạn. Dù bộ
kiểm thử thực tế không thiết kế dữ liệu nhắm vào trường hợp này để giữ sự đơn
giản, xét về tính chặt chẽ của thuật toán ta vẫn phải coi trọng vấn đề đó.
Để tránh trường hợp này, có thể dừng chia đôi khi hiệu hai cận nhỏ hơn một số
dương nhỏ, rồi lấy kết quả làm tròn lên từ cận trên.

Đến đây thuật toán đã hoàn chỉnh.

\paragraph{Tiểu kết.}
Giải các bài toán động thường là một điểm khó trong hình học tính toán. Với
mọi bài thuộc kiểu này, khi lập trình cần cân nhắc cẩn thận nhiều trường hợp,
như phạm vi giá trị của biến, cách lấy hay bỏ biên, và cách tránh để các phép
tính phức tạp khuếch đại sai số dấu phẩy động.

Bài này cũng có thể giải thuận lợi bằng phương pháp truyền thống, nhưng tính
toán vất vả và quy trình phức tạp. Nhờ dùng chia đôi cùng tính đơn điệu ẩn
trong đề bài, ta biến động thành tĩnh, giảm đáng kể độ phức tạp tư duy và độ
phức tạp lập trình.

\subsection{Ví dụ 3: \textit{Heliport} (2004 ACM/ICPC World Finals)}

\paragraph{Ý nghĩa bài toán.}
Cho mái nhà là một đa giác đơn \(N\) cạnh \((1\le N\le 20)\), chỉ gồm các
cạnh ngang và dọc; mỗi cạnh có độ dài nguyên dương không quá \(50\). Cần xây
một sân bay trực thăng hình tròn trên mái nhà. Hãy tìm bán kính lớn nhất có
thể.

\paragraph{Phân tích.}
Mô tả của bài toán khá rõ: tìm đường tròn nội tiếp lớn nhất trong một đa giác
đơn trực giao \(P\). Giả sử đoạn ngang \(i\) có tung độ \(h_i\), hai đầu mút
có hoành độ \(l_i,r_i\) với \(l_i<r_i\); đoạn dọc \(j\) có hoành độ \(v_j\),
hai đầu mút có tung độ \(d_j,u_j\) với \(d_j<u_j\). Một mô hình toán học tự
nhiên là
\[
\begin{array}{ll}
\text{cực đại hóa} & r,\\[0.2em]
\text{với điều kiện} & (x,y)\in P,\quad r\ge 0,\\
& \dist\bigl((x,y),e\bigr)\ge r
  \quad\text{với mọi cạnh biên } e \text{ của } P.
\end{array}
\]
Khi khai triển khoảng cách tới các đỉnh và tới các cạnh ngang, dọc, đây là
một bài toán quy hoạch bậc hai. Bằng cách làm cho ba đẳng thức trong các điều
kiện trở thành dấu bằng, rồi kiểm tra các điều kiện còn lại, ta có thể thu
được một nghiệm; nghiệm tối ưu chắc chắn nằm trong số đó.

Nhờ mô hình trên, ta đã đại số hóa thành công một bài toán hình học. Vì
\(N\) trong đề rất nhỏ, không cần dùng các thuật toán chuyên dụng phức tạp cho
quy hoạch bậc hai; chỉ cần liệt kê, cho ba đẳng thức thành lập rồi kiểm chứng
các quan hệ khác. Các công thức trong bài có ba loại: khoảng cách tới điểm
(ký hiệu \(P\)), khoảng cách tới cạnh ngang (ký hiệu \(H\)), và khoảng cách
tới cạnh dọc (ký hiệu \(V\)). Ba đẳng thức cùng thành lập có mười trường hợp:
\[
  PPP,\ PPH,\ PPV,\ PHH,\ PVV,\ PHV,\ HHV,\ HVV,\ HHH,\ VVV.
\]

Xét riêng từng trường hợp: trước hết \(HHH\) và \(VVV\) là không thể, vì ba
đường thẳng song song không xác định được một đường tròn. Tám trường hợp còn
lại \(PPP,PPH,PPV,PHV,HHV,HVV,PHH,PVV\) đều không thể bỏ qua.

Bằng cách lập phương trình cho từng trường hợp, ta có thể giải ra tọa độ tâm
đường tròn, và bài toán tưởng như đã được giải triệt để. Nhưng trong thao tác
thực tế, không khó nhận ra việc tính phương trình rất phức tạp, lại có khả
năng mẫu số bằng \(0\) hoặc xuất hiện nghiệm phức. Ngay cả khi dùng phần mềm
để giải, lúc lập trình vẫn rất dễ sai do nhập không cẩn thận hoặc xét thiếu
trường hợp, kéo theo rắc rối khi gỡ lỗi và làm tăng độ khó để qua kiểm thử.
Vì vậy tốt nhất ta đổi cách nghĩ để giảm lượng tính toán.

Ta vừa thấy tám trường hợp trên đã là tối ưu. Chẳng lẽ không còn cách nào
tiếp tục đơn giản hóa? Trước hết xét điều kiện xác định hữu hạn số đường
tròn: ba khoảng cách bằng nhau và bằng bán kính, mỗi khoảng cách tương ứng
với một điểm biên hoặc một tiếp tuyến. Nếu biết bán kính, ta chỉ cần thêm hai
khoảng cách bằng bán kính là có thể xác định đường tròn; giảm đi một công
thức khoảng cách cũng vẫn xác định được nghiệm.

Tám trường hợp nói trên lần lượt là \(PPP,PPH,PPV,PHV,HHV,HVV,PHH,PVV\).
Với mỗi trường hợp, nếu thay một khoảng cách bằng bán kính, ta rút gọn còn
bốn loại:
\[
  PP,\ PH,\ PV,\ HV.
\]
Như vậy lượng tính toán giảm mạnh, việc giải phương trình đơn giản hơn nhiều,
và xác suất phát sinh lỗi cũng thấp đi.

Nhưng thứ bài toán yêu cầu lại chính là bán kính; bán kính là ẩn số, không
thể coi là đại lượng đã biết để lập công thức. Đúng lúc ta tưởng như bế tắc,
tư tưởng chia đôi lại cho biết bí quyết giải bài: chia đôi \(r\).

Tính đơn điệu của bài toán rất rõ ràng. Về mặt hình học, nghiệm tối ưu nếu
thu nhỏ tại vị trí cũ vẫn hợp lệ, còn phóng to sẽ vượt qua biên. Về mặt đại
số, các điều kiện ràng buộc liên quan tới bán kính đều cùng chiều, đều là
``không nhỏ hơn khoảng cách yêu cầu'' hoặc tương đương.

Sau khi chia đôi \(r\), ta lập phương trình theo bốn loại \(PP,PH,PV,HV\) để
kiểm tra tính hợp lệ của nghiệm. Dựa vào kết quả kiểm tra, liên tục điều chỉnh
cận cho đến khi đạt yêu cầu độ chính xác, ta thu được đáp án.

\paragraph{Tiểu kết.}
Hiểu theo mặt chữ, ``hình học tính toán'' rõ ràng không thể tách khỏi
``tính toán''. Tuy nhiều bài nhìn qua không quá khó, phần tính toán lại khá
phức tạp; nếu không chú ý chi tiết, rất có thể sinh ra những lỗi nhỏ ngoài dự
kiến, khiến đáp án ``sai một ly, đi một dặm''. Trong một số tình huống, tư
tưởng chia đôi có thể giảm mạnh lượng tính toán, đồng thời hạ thấp độ phức
tạp tư duy và độ phức tạp lập trình.

Bài này thuộc kiểu ``dễ mà rườm rà'': bản thân thuật toán không khó, nhưng
tính toán và lập trình rất phiền. Với bài này, chữ ``phân'' trong chia đôi
không phá vỡ cấu trúc; ngược lại, nó gom những trường hợp rời rạc lại với
nhau, biến lẻ tẻ thành chỉnh thể và qua đó tối ưu hóa thuật toán.

\subsection{Ví dụ 4: \textit{Flight Safety} (2007 ACM/ICPC NWERC)}

\paragraph{Ý nghĩa bài toán.}
Có \(C\) lục địa đôi một không giao nhau \((1\le C\le 20)\), mỗi lục địa là
một đa giác đơn có số cạnh \(M\) \((3\le M\le 30)\). Đường bay là một đường
gấp khúc gồm \(N\) đoạn thẳng \((2\le N\le 20)\). Hãy tìm giá trị lớn nhất,
trên các điểm của đường bay, của khoảng cách từ điểm đó tới lục địa gần nhất.

\paragraph{Phân tích.}
Đây là một bài toán khoảng cách hình học. Vì đường bay gồm nhiều đoạn thẳng,
ta có thể xét riêng từng đoạn; bài toán được quy về tìm giá trị lớn nhất của
khoảng cách từ một đoạn thẳng tới nhiều đa giác.

Nếu trực tiếp giải khoảng cách giữa một đoạn thẳng và nhiều đa giác thì khá
phiền. Do hai hoặc nhiều đa giác tác động qua lại, điểm cực trị không nhất
thiết là đầu mút của đoạn bay hay hình chiếu của một đỉnh đa giác lên đoạn
bay. Muốn tìm cực trị phải xét quan hệ tương hỗ giữa các đa giác, đây là một
vấn đề khó.

Để tìm lời giải, hãy giả sử ta đã biết một đáp án và hỏi làm thế nào để kiểm
tra tính hợp lệ và tối ưu của nó. Nếu đáp án hợp lệ, mọi điểm trên đường bay
đều có khoảng cách tới các đa giác nhỏ hơn hoặc bằng đáp án ấy. Nhìn từ góc
độ khác, điều này tương đương với việc mọi điểm trên đường bay đều nằm trong
các đa giác đã được ``nở'' ra ngoài thêm đúng khoảng cách đó. Đa giác sau khi
``nở'' trông như thế nào? Với đỉnh, nó nở thành một cung tròn; với cạnh, nó
nở ra ngoài thành một hình chữ nhật có bề rộng tương ứng.

Nếu một đỉnh của đường bay nằm trên biên của đa giác đã nở, ta vừa chạm tới
đáp án tối ưu. Nếu một phần của đường bay nằm ở bên ngoài, đáp án không hợp
lệ. Nếu đường bay không chạm biên, đáp án chưa tối ưu.

Để phán định quan hệ trong--ngoài, ta có thể tìm các giao điểm giữa từng đoạn
bay với từng cạnh hoặc từng đường tròn của mỗi lục địa. Cắt đoạn bay theo các
giao điểm ấy, mỗi đoạn con sẽ hoặc nằm trọn trong một hình, hoặc nằm trọn
ngoài hình đó. Khi ấy trạng thái trong--ngoài của đoạn con tương đương với
trạng thái của trung điểm: chỉ cần trung điểm nằm trong một hình tròn, một
hình chữ nhật, hoặc trong đa giác gốc, thì đoạn con nằm trong hình; ngược lại
nó ở ngoài, và đáp án không hợp lệ.

Rõ ràng khoảng cách càng lớn thì phần nở ra càng lớn, nên bài toán có tính
đơn điệu. Muốn tìm nghiệm tối ưu, chia đôi lại giúp ta: nếu còn điểm bên
ngoài thì tăng đáp án, ngược lại giảm đáp án, cho đến khi hai cận đủ gần rồi
in kết quả.

\paragraph{Tiểu kết.}
Những bài có nhiều phần tử hình học và quan hệ tương hỗ phức tạp như bài này
rất thường gặp trong hình học tính toán. Giải trực tiếp thường khó, việc mã
hóa và gỡ lỗi cũng gặp nhiều rắc rối. Tư tưởng chia đôi có thể loại bỏ các
ràng buộc không cần thiết, chỉ giữ lại phần then chốt, nhờ đó làm bài toán
đơn giản hơn rất nhiều.

Trong bài này, các đa giác ràng buộc lẫn nhau nên khó tìm cực trị. Phương
pháp chia đôi tránh được điểm khó đó, chuyển bài toán tối ưu hóa thành bài
toán phán định, biến ``tìm'' thành ``chứng'', khiến lập trình nhẹ nhàng hơn,
thuật toán giản dị hơn, và vấn đề được giải quyết thuận lợi.

\section{Tổng Kết}

Hình học tính toán rộng lớn và sâu sắc; các bài toán liên quan biến hóa muôn
hình vạn trạng, cách giải cũng không có khuôn mẫu cố định. Một vài bài kinh
điển chỉ cần sửa đổi đôi chút là phương pháp truyền thống đã không còn lợi
thế. Điều đó đòi hỏi ta phải thoát khỏi quán tính tư duy và dùng những ý
tưởng mới; tư tưởng chia đôi chính là một thành viên không thể thiếu trong số
đó.

Qua phân tích các ví dụ trên, ta có thêm nhận thức mới về ứng dụng của tư
tưởng chia đôi trong hình học tính toán. Có được tư tưởng ấy, ta có thể biến
phức tạp thành đơn giản, biến động thành tĩnh, biến lẻ tẻ thành chỉnh thể,
biến tìm kiếm thành chứng minh; dùng cách đơn giản để giải bài, giảm bớt xử
lý chi tiết rườm rà và quan hệ ràng buộc phức tạp, từ đó thu được kết quả
thỏa đáng một cách gọn gàng, hiệu quả. Có thể nói, chia đôi đem lại cho ta
một cách nghĩ hoàn toàn mới, là công cụ sắc bén để giải thành công các bài
toán hình học tính toán.

\section*{Lời Cảm Ơn}

Xin cảm ơn thầy Liu Rujia đã cung cấp ví dụ 3 và ví dụ 4 cho bài viết này.

Xin cảm ơn các thầy cô và các bạn học sinh đã góp nhiều ý kiến, đề xuất cho
bài viết.

\section*{Tài Liệu Tham Khảo}

\begin{enumerate}
  \item Liu Rujia, Huang Liang. \textit{Nghệ thuật thuật toán và thi tin học}.
  Bắc Kinh: Nhà xuất bản Đại học Thanh Hoa, 2004.
  \item Zhou Peide. \textit{Hình học tính toán: thiết kế và phân tích thuật
  toán}, ấn bản thứ 2. Bắc Kinh: Nhà xuất bản Đại học Thanh Hoa, 2005.
  \item Mark de Berg, Marc van Kreveld, Mark Overmars, Otfried Schwarzkopf.
  \textit{Computational Geometry: Algorithms and Applications}, 2nd Edition.
  Springer, 2000.
\end{enumerate}

\section*{Phụ Lục}

\noindent Bài gốc của ví dụ 1:
\url{http://acmicpc-live-archive.uva.es/nuevoportal/data/problem.php?p=3270}

\noindent Bài gốc của ví dụ 2:
\url{http://acmicpc-live-archive.uva.es/nuevoportal/data/problem.php?p=2397}

\noindent Bài gốc của ví dụ 3:
\url{http://acmicpc-live-archive.uva.es/nuevoportal/data/problem.php?p=2994}

\noindent Bài gốc, dữ liệu và lời giải của ví dụ 4:
\url{http://2007.nwerc.eu/problems/nwerc07-problemset.pdf}, Problem F.

\noindent Dữ liệu:
\url{http://2007.nwerc.eu/problems/nwerc07-contest-testdata.zip}

\noindent Lời giải:
\url{http://2007.nwerc.eu/problems/nwerc07-solutions.zip}.

\end{document}
